% 03_benchmark -- front lane.
% src (counts): juliet/full_split_v1.json; juliet_full/data/README.md; juliet_full/data/build/verify_report.md;
%      juliet_full/data/build/verify_layer3_report.md; juliet_full/lane_reports/data_stage3.md; probes via \res{x}{probe_*}.
\section{Benchmark Construction}
\label{sec:benchmark}

\subsection{Source}
We start from all \cwe{} test cases of Juliet C/C++ v1.3~\cite{black2018juliet13,sardjuliet13}.
They form 1,658 groups from 42 families and 48 flow variants, 868 in C and 790 in C++ (Table~\ref{tab:dataset}). % src: juliet/full_split_v1.json
Each group yields one bad and one good \emph{unit}: the source files that a harness compiles, links and runs.
We tag every group with its leak mechanism.
\code{c\_alloc} (960 groups) leaks a block from \code{malloc}, \code{calloc}, \code{realloc} or \code{strdup}.
\code{cpp\_new} (588) leaks an object or array from \code{new}.
\code{realloc\_fail} (108) leaks the old block when \code{realloc} fails.
\code{destructor} (2) leaks through a class without a proper destructor.

\subsection{Shortcut audit}
\label{sec:audit}
Before we used Juliet, we asked whether a detector could score well without reading the code.
We ran six probes on the eval units (Table~\ref{tab:probes}); thresholds and the bag-of-words (BoW) logistic regression are fitted on dev.
On \Orig, the presence of a \code{static void} definition separates the labels perfectly (\res{x}{probe_static_void_original}).
In Juliet the good functions are \code{static} and are called from a non-static \code{good()} wrapper; every good file contains a \code{static void} definition and no bad file does. % src: juliet_full/data/build/probes_report.md (good 100.0%, bad 0.0%)
Function count (\res{x}{probe_func_count_original}) and length (\res{x}{probe_length_lines_original}) work because a good file holds two good functions and a wrapper.
The grep rule ``no \code{free}/\code{delete} $\Rightarrow$ leak'' reaches \res{x}{probe_grep_v0_original}, because the B2G function in every good file contains the fixing \code{free}.
None of these features is about leaks.
For \Deb{} we fixed an acceptance limit before the build: macro-F1 $\le 0.6$ for every format probe. % src: juliet_full/data/README.md (probe acceptance limit)
On \Deb{} the format probes fall to \res{x}{probe_static_void_layer1}--\res{x}{probe_length_lines_layer1}, and grep to \res{x}{probe_grep_v0_layer1}, because G2B good units use stack memory and have no \code{free} either.
A BoW model still reaches \res{x}{probe_bow_lr_layer1}, and this cannot be removed without changing semantics: a G2B good unit allocates on the stack where the bad unit calls \code{malloc} or \code{new}, and a B2G good unit contains a \code{free} that the bad unit lacks.
\Dead{} (\S\ref{sec:dead}) removes the second cue.

\IfFileExists{tables/res_probes.tex}{% generated by figures/make_figures.py -- do not edit
% src: juliet_full/analysis/score.py on juliet_full/data/build/probes_layer3.json (eval set)
\begin{table}[t]
\centering
\caption{Shortcut probes (macro-F1 on the 2,656 eval units; thresholds and the bag-of-words model are fitted on dev). The last row is the bag-of-words AUROC. A one-class predictor scores 0.333.}
\label{tab:probes}
\scriptsize
\setlength{\tabcolsep}{3pt}
\begin{tabular}{@{}lrrrr@{}}
\toprule
Probe & \Orig{} & \Deb{} & \Dead{} & \Plac{} \\
\midrule
\code{static void} $\Rightarrow$ good & 1.000 & 0.333 & 0.333 & 0.333 \\
function count & 0.924 & 0.499 & 0.499 & 0.499 \\
length (lines) & 0.906 & 0.544 & 0.525 & 0.498 \\
length (tokens) & 0.851 & 0.501 & 0.511 & 0.520 \\
grep: no \code{free}/\code{delete} $\Rightarrow$ leak & 0.935 & 0.701 & 0.333 & 0.701 \\
grep v1: \code{free} outside \code{if} $\Rightarrow$ good & 0.935 & 0.656 & 0.333 & 0.668 \\
grep v2: free-count threshold & 0.935 & 0.701 & 0.498 & 0.701 \\
bag-of-words LR & 1.000 & 0.979 & 0.752 & 0.981 \\
bag-of-words AUROC & 1.000 & 0.982 & 0.810 & 0.988 \\
\bottomrule
\end{tabular}
\end{table}
}{}
% front lane. src: juliet/full_split_v1.json (groups[], summary) -- counts recomputed from groups[]
\begin{table}[t]
\centering
\caption{Juliet v1.3 \cwe{} groups by flow category and leak mechanism. Each group gives one bad and one good unit in every layer.}
\label{tab:dataset}
\footnotesize
\setlength{\tabcolsep}{4pt}
\begin{tabular}{@{}lrrrrrr@{}}
\toprule
Flow category & variants & groups & C & C++ & dev & eval \\
\midrule
01 baseline & 1 & 42 & 26 & 16 & 0 & 42 \\
02--22 control flow & 19 & 748 & 482 & 266 & 160 & 588 \\
31--45 data flow & 9 & 278 & 140 & 138 & 68 & 210 \\
51--68 across files & 12 & 380 & 220 & 160 & 68 & 312 \\
71--84 C++ objects & 7 & 210 & 0 & 210 & 34 & 176 \\
\midrule
Mechanism & families & & & & & \\
\midrule
\code{c\_alloc} (\code{malloc} \dots) & 20 & 960 & 760 & 200 & 180 & 780 \\
\code{cpp\_new} (\code{new}, \code{new[]}) & 14 & 588 & 0 & 588 & 126 & 462 \\
\code{realloc\_fail} & 6 & 108 & 108 & 0 & 24 & 84 \\
\code{destructor} & 2 & 2 & 0 & 2 & 0 & 2 \\
\midrule
Total & 48 / 42 & 1,658 & 868 & 790 & 330 & 1,328 \\
\bottomrule
\end{tabular}
\end{table}


\subsection{Construction of \Deb}
\label{sec:deb}
\noindent\textbf{Good entry.}
Each good file contains a G2B and a B2G function; we keep exactly one, so that a good unit has the shape of a bad unit.
The choice is stratified by set, flow category and mechanism: half of each stratum gets G2B and half B2G, and an odd stratum is decided by a seeded coin flip (seed 20261123). % src: juliet_full/data/README.md
The result is 621 G2B and 621 B2G good units in eval. % src: \res{x}{n_kind_G2B}, \res{x}{n_kind_B2G}
The 110 \code{realloc\_fail} and \code{destructor} groups have numbered good functions (\code{good1}, \code{good2}) instead of a G2B/B2G pair; we pick one with a per-group seeded generator. % src: juliet_full/data/README.md (good_kind other = 110)
Only code reachable from the chosen entry is kept; the bad unit is the \code{OMITGOOD} view of the file.

\noindent\textbf{Format.}
We remove \code{static} from every function definition and prototype in both labels; file-scope variables keep it.
The entry becomes \code{void sym\_0()}, called by the same driver \code{main} in both labels.
Comments are removed, the \code{OMITBAD}/\code{OMITGOOD}/\code{INCLUDEMAIN} macros are resolved, and every Juliet identifier is renamed to \code{sym\_N}.
A leakage test checks that no identifier, comment or string in any of the 6,632 records across \Deb{} and \Orig{} matches \code{bad}, \code{good}, \code{g2b}, \code{b2g}, \code{cwe}, \code{leak}, \code{omit}, \code{flaw} or \code{fix}. % src: juliet_full/data/README.md (leakage test passes all 6,632 records)

\subsection{Label verification}
\label{sec:verify}
All 3,316 \Deb{} units compile and link with both GCC~13 and Clang~23. % src: juliet_full/data/build/verify_report.md
We run every unit under ASan/LSan (\code{-O0 -g}, \code{detect\_leaks=1}, 10\,s timeout, stdin closed).
We set \code{use\_stacks=0}: at \code{-O0} stale stack slots keep leaked blocks ``reachable'', and the default setting misses 175 more bad units. % src: verify_report.md
No good unit leaks (1,658/1,658).
For bad units, LSan reports a leak in 1,474 of 1,658 cases (\code{c\_alloc} 914/960, \code{cpp\_new} 558/588, \code{destructor} 2/2). % src: verify_report.md
All 184 misses have a documented cause, none a labeling error: 68 units (variants 45, 68) leak a block still held by a global at exit, 108 \code{realloc\_fail} units leak only if \code{realloc} fails, and 8 units of variant~12 take the leaking branch only on some runs (\code{rand()}; we run up to six times).
Valgrind~3.28 on a stratified sample of 200 units agrees: it reports no leak for 100/100 good units, and for the 100 bad units it agrees with LSan on 99 and with the label on 87. % src: verify_report.md
The dynamic-enhanced reference (\S\ref{sec:methods}) closes the three documented gaps by fault injection.

\subsection{\Dead: never-executed frees}
\label{sec:dead}
A detector that reasons about paths must ignore a \code{free} that cannot run.
\Dead{} tests this by making every unit contain the same kind of release statements, live or dead.
\begin{itemize}
\item \emph{Flaw sites.} Where the good sink frees and the bad sink does not, we insert a dead release of the sink pointer under an always-false guard. G2B good units, whose sink equals the bad sink, get the same dead release.
\item \emph{Real frees.} Every real \code{free}/\code{delete}/\code{delete[]} in both labels is wrapped in an always-true guard, so a live release also sits inside an \code{if}.
\item \emph{Matching forms.} The dead form matches the allocation (\code{free}, \code{delete}, \code{delete[]}); empty destructors get a dead release inside. For \code{realloc}, the dead free of the new pointer follows the call.
\item \emph{Count matching.} Each group gets $e\in\{0,1\}$ extra dead releases at the same anchor in both labels, chosen by a seeded generator, so the number of release statements and guards is equal within every bad/good pair.
\end{itemize}
Guards compare a file-local constant \code{static const int sym\_N = V;} ($V\in\{4,5\}$, $N$ random in 10--99) in six forms (e.g.\ \code{==5}, \code{>4}, \code{!=4}), each used as both true and false guard. % src: juliet_full/data/README.md (layer3 guards)
We chose this on dev with a leave-variant-out BoW probe on B2G pairs: a fixed constant gives AUROC 0.989, our design 0.572. % src: juliet_full/lane_reports/data_stage3.md (dev design check)
On eval, every grep variant falls to $\le$\res{x}{probe_grep_v2_layer3}, and the BoW AUROC is \res{x}{probe_bow_auroc_layer3}.
The remaining BoW signal comes from G2B units, whose stack allocation is intrinsic to the label. % src: juliet_full/data/build/probes_layer3_report.md (B2G .609, G2B .986)

All 3,316 \Dead{} units compile with GCC and Clang; all 3,236 units outside variant~12 keep their \Deb{} dynamic verdict. % src: juliet_full/data/build/verify_layer3_report.md

\subsection{\Plac{} and \Orig}
\label{sec:plac}
\Plac{} uses the same insertion points, guards and counts as \Dead, but the dead statement is a print call, \code{printLine("Status ok")} or one of three similar strings.
Real frees are not wrapped, so a drop that also appears on \Plac{} is an editing effect.
\Plac{} passes the same compile and dynamic checks. % src: verify_layer3_report.md (placebo 3316/3316, same verdicts outside variant 12)
\Orig{} pairs the \Deb{} bad unit with the complete original good program (both static good functions and the wrapper).

\subsection{Split and subsets}
\label{sec:split}
We split by flow variant, not by family (Table~\ref{tab:split}).
Families differ only in data type and API, so a family split would leave each control-flow template on both sides.
The dev set has nine variants (04, 10, 17, 18, 31, 41, 65, 68, 73), drawn per flow category with quotas 4/2/2/1 from variants that contain \code{cpp\_new} (Appendix~\ref{app:split}). % src: juliet/full_split_v1.json (rules)
This gives 330 dev groups (660 units) and 1,328 eval groups (2,656 units) over 39 variants.
Trained methods use three folds of eval variants: fold $k$ is predicted by a model trained on the other two, so no eval unit is scored by a model that saw its variant.
Nested subsets (stratified by flow category and mechanism) serve the budget-limited DeepSeek runs, repeats and latency.

% front lane. src: juliet/full_split_v1.json (dev_variants, fold_variants, summary); subsets: juliet_full/data/build/subsets.json
\begin{table}[t]
\centering
\caption{Flow-variant split and evaluation subsets. All test results use eval units. Trained models use 3-fold grouped cross-validation over eval (train on two folds, test on the third) and dev only for validation. Variant lists are in Appendix~\ref{app:split}.}
\label{tab:split}
\footnotesize
\setlength{\tabcolsep}{3pt}
\begin{tabularx}{\columnwidth}{@{}lrrr>{\raggedright\arraybackslash}X@{}}
\toprule
Set & variants & groups & units & use \\
\midrule
dev & 9 & 330 & 660 & tuning, selection \\
eval & 39 & 1,328 & 2,656 & all reported scores \\
\quad fold 0 & 15 & 472 & 944 & CV test fold \\
\quad fold 1 & 12 & 424 & 848 & CV test fold \\
\quad fold 2 & 12 & 432 & 864 & CV test fold \\
\midrule
\code{api\_subset} & -- & 664 & 1,328 & DeepSeek L1--L3 on \Dead \\
\code{nested25} & -- & 332 & 664 & DeepSeek paired \Dead/\Plac/\Orig \\
\code{repeat100} & -- & 100 & 200 & LLM repeat runs \\
\code{latency300} & -- & 150 & 300 & latency (concurrency 1) \\
\bottomrule
\end{tabularx}
\end{table}

